\documentclass[10pt,a4paper]{amsart}
\usepackage[margin=19mm]{geometry}
\usepackage{amsmath,amssymb,mathtools}
\usepackage[scr=rsfs,cal=euler]{mathalfa}
\usepackage{xfrac}
\usepackage[unicode,hidelinks,bookmarks=false]{hyperref}
\newcommand{\ie}{i.e.\ }
\newcommand{\cf}{cf.\ }
\newcommand{\resp}{respectively\ }
\newcommand{\Subsection}{\subsection}
\providecommand{\nobreakdash}{\nobreak\hbox{-}}
\setlength{\emergencystretch}{2em}
\multlinegap0pt
\usepackage{bm}
\usepackage{bbm}
\usepackage{dsfont}
\usepackage{xspace}
\usepackage{cancel}
\usepackage{mathtools}
\usepackage{amsthm}
\makeatletter
\@ifundefined{lem}{\newtheorem{lem}{Lemma}}{}
\makeatother
\title[Absence of embedded spectrum for nonlinear Schr\"odinger equ]{Absence of embedded spectrum for nonlinear Schr\"odinger equations linearized around one dimensional ground states}
\author{Charles Collot, Pierre Germain, Eliot Pacherie}
\date{Source: arXiv:2503.02957v2 (CC BY 4.0)}
\begin{document}
\maketitle

\noindent\emph{Source and licence.} Absence of embedded spectrum for nonlinear Schr\"odinger equations linearized around one dimensional ground states. Version arXiv:2503.02957v2 (CC BY 4.0), distributed under the Creative Commons Attribution 4.0 International licence, as recorded in the arXiv licence metadata for this exact version and shown on the abstract page https://arxiv.org/abs/2503.02957v2. Full rights notice: licenses/arxiv-2503.02957-CC-BY-4.0.txt.\par\smallskip

\noindent\textit{Editorial scope.} Counted unit: the Lem whose source label is \texttt{fiberoptic} in the article (source0.tex lines 211--239, source label \texttt{fiberoptic}), delivered together with the article's own complete original proof of it (source0.tex lines 241--325, one \texttt{proof} environment of 85 lines). No mathematics is written, repaired or re-derived: statement and proof are the article's own text line for line. Required context retained uncounted: kanine (lines 191--195); avb (lines 199--203). Every definition, display and label of those lines is retained unchanged. Omitted from this package and not redistributed: 1--190, 196--198, 204--210, 240--240, 326--467. Nothing inside a retained excerpt is omitted or altered: every retained body is delivered exactly as the article's lines are written, so no omission receipt of any kind accompanies this package. Citations: the retained excerpts carry no citation command at all, so no bibliography entry of the article is transcribed and no reference is redistributed. Reference adaptation: none was needed and none was made; every reference inside the delivered unit resolves inside it, so no unresolved reference or citation is produced. Printed slips kept literal: no printed word, symbol, hypothesis or number of the counted statement or of its proof was altered. Layout and class: the article's own class is \texttt{amsart} with packages including inputenc, tikz, tikz-3dplot, amsmath, amsthm, amssymb, graphics, esint, amscd, color, xcolor, mathtools, url, bbm; the wrapper is the lane's pinned \texttt{amsart} editorial wrapper at 10pt on A4, which restates the theorem-like environments the retained text needs and adds the article's own macro definitions that the retained text needs (none), with extra wrapper packages (bm, bbm, dsfont, xspace, cancel, mathtools). The delivered numbering is the unit's own. Assets: no retained span contains a figure environment, a graphics-inclusion command, a PDF-page inclusion, a table or a TikZ picture; no image file is copied into this package, the package declares no asset dependency and no image file is redistributed with it. Licence: version arXiv:2503.02957v2 of this article is distributed under the Creative Commons Attribution 4.0 International licence, as recorded in the arXiv licence metadata for this exact version and shown on the abstract page https://arxiv.org/abs/2503.02957v2. A full rights notice is delivered at licenses/arxiv-2503.02957-CC-BY-4.0.txt. Characters: every retained excerpt is pure ASCII, so the delivered engine, XeLaTeX with its default Latin Modern text font, reports no missing character.
\begin{equation}
  Q (x) = c_0 e^{-\sqrt{\mu} x} (1 + o_{x \rightarrow
  + \infty} (1)),Q'(x)=-\sqrt{\mu}c_0 e^{-\sqrt{\mu} x} (1+o_{x \rightarrow
  + \infty} (1)). \label{kanine}
\end{equation}

\begin{equation}
  f (x) = e^{- \sqrt{\mu} x} - e^{\sqrt{\mu} x} \int_x^{+ \infty} e^{- 2
  \sqrt{\mu} t} \int_t^{+ \infty} F (Q^2) (s) e^{\sqrt{\mu} s} f (s) d s d t
  \label{avb}
\end{equation}

\begin{lem}
  \label{fiberoptic}Given $\lambda \in  \sigma_e (\mathcal{H})$, any nonzero solution
  of
  \[ (\mathcal{H}- \lambda) \left(\begin{array}{c}
       f\\
       g
     \end{array}\right) = 0 \]
  that is in $L^2 (\mathbb{R}, \mathbb{R})$ must satisfy, for some $c_* \neq 0$,

\begin{align*}
& \left(\begin{array}{c}
       f\\
       g
     \end{array}\right) = c_{\ast} e^{-\sqrt{\mu + \lambda} x} \left( \left(\begin{array}{c}
       1\\
       0
     \end{array}\right) + o_{x \rightarrow + \infty} (1) \right)
      \qquad \mbox{if $\lambda \geq \mu$} \\
& \left(\begin{array}{c}
       f\\
       g
     \end{array}\right) = c_{\ast} e^{-\sqrt{\mu - \lambda} x}
     \left( \left(\begin{array}{c}
       0\\
       1
     \end{array}\right) + o_{x \rightarrow + \infty} (1) \right) \qquad \mbox{if $\lambda \leq - \mu$}.
\end{align*}

\end{lem}

\begin{proof}
  We recall that the equation $\mathcal{H}- \lambda = 0$ is
  \[ \left(\begin{array}{cc}
       \partial_x^2 - \mu - \lambda & 0\\
       0 & - (\partial_x^2 - \mu + \lambda)
     \end{array}\right) \left(\begin{array}{c}
       f\\
       g
     \end{array}\right) = - V \left(\begin{array}{c}
       f\\
       g
     \end{array}\right) . \]
  
  \
  
\noindent  1) If $\lambda > \mu > 0$, if we consider the equation with $V = 0$, the
  radial solutions of $\partial_x^2 - \mu - \lambda = 0$ are $e^{-\sqrt{\mu + \lambda} x} $ and $e^{\sqrt{\mu + \lambda} x} $, while the solutions of $\partial_x^2 - \mu +
  \lambda$ are $\cos \left( \sqrt{\lambda - \mu } x \right)$ and
  $\sin \left( \sqrt{\lambda - \mu } x \right)$.
  
  Since $V$ decays exponentially fast by (\ref{kanine}), we construct by a
  fixed point argument four solutions of $\mathcal{H}- \lambda = 0$ near $x =
  + \infty$ which satisfy respectively

  \begin{eqnarray*}
    \left(\begin{array}{c}
      f_1\\
      g_1
    \end{array}\right) & = & e^{-\sqrt{\mu + \lambda } x}\left( \left(\begin{array}{c}
       1\\
       0
     \end{array}\right) + o_{x \rightarrow + \infty} (1) \right)\\
    \left(\begin{array}{c}
      f_2\\
      g_2
    \end{array}\right) & = & e^{\sqrt{\mu +  \lambda} x} \left( \left(\begin{array}{c}
       1\\
       0
     \end{array}\right) + o_{x \rightarrow + \infty} (1) \right)\\
    \left(\begin{array}{c}
      f_3\\
      g_3
    \end{array}\right) & = & \cos \left( \sqrt{ \lambda - \mu }\, x
    \right)\left( \left(\begin{array}{c}
       0\\
       1
     \end{array}\right) + o_{x \rightarrow + \infty} (1) \right)\\
    \left(\begin{array}{c}
      f_4\\
      g_4
    \end{array}\right) & = & \sin \left( \sqrt{ \lambda - \mu }\, x
    \right)\left( \left(\begin{array}{c}
       0\\
       1
     \end{array}\right) + o_{x \rightarrow + \infty} (1) \right) .
  \end{eqnarray*}
  They can indeed be constructed by a fixed point similar to (\ref{avb}), replacing $F(Q^2)$ by $V$, $\mu$ by $\mu+\lambda$ or $\mu-\lambda$ (with the possibility for their square root to be imaginary) and for vector valued functions, using the decay of $V$ to close the fixed point argument for $r$ large enough.

\
  
  Since $\mathcal{H}- \lambda = 0$ is a fourth order ODE,
  any solution must be a linear combination of these functions. Looking at
  their behavior near $x = + \infty$, the only possibility for such a
  combination to be in $L^2 (\mathbb{R}, \mathbb{R})$ is for
  the solution to be colinear to $\left(\begin{array}{c}
    f_1\\
    g_1
  \end{array}\right)$.
  
  \
  
\noindent  2)  If $\lambda < - \mu$, the solutions of $\partial_x^2 - \mu - \lambda$ are then
  $\cos \left( \sqrt{| \lambda + \mu |} x \right)$ and $\sin
  \left( \sqrt{| \mu + \lambda |} x \right)$, while the solutions of $\partial_x^2 -
  \mu + \lambda$ are $e^{-\sqrt{\mu -  \lambda } x}$
  and $e^{\sqrt{\mu -  \lambda } x}$. We conclude
  as in the first case, but now the exponentially decaying function is at the limit on the
  vector $\left(\begin{array}{c}
    0\\
    1
  \end{array}\right)$.

\noindent 3) If $\lambda = \mu$ or $\lambda = - \mu$, instead of the two oscillating functions, it is $1$ and $x$ that are solutions for $V=0$. We can still construct two solutions of $\mathcal{H}-\lambda$ with these behaviors near $x = +\infty$ by the same fixed point argument, and any linear combination of them do not belong in $L^2$, leading to the same conclusion.
  
\end{proof}

\end{document}
